\chapter{Conclusion}
\label{chap:conclusion}
This thesis investigates how numerical variability propagates through resting-state fMRI functional connectivity pipelines into graph-theoretical measures and downstream statistical inference. Phase~I demonstrates that numerical variability in fMRIPrep-processed PPMI networks is measurable relative to inter-subject variability, depends on graph metric, brain region, threshold, population, and confound-regression choices, and can induce significance flips for results near decision thresholds across sample sizes. Phase~II builds on these findings to test whether numerical varaibility can be estimated from a single standard connectome without repeated Monte Carlo Arithmetic executions. Preliminary results indicate that compact features capturing threshold proximity and edge-toggle influence provide strong predictive signal across many metric--threshold combinations, motivating systematic model evaluation, cross-cohort and external validation, and lightweight approximation of NPVR and significance-flip risk to support more numerically reliable network neuroscience.